%% ----------------------------------------------------------------------------
%%  §5.7  Appendice: le funzioni trigonometriche inverse
%% ----------------------------------------------------------------------------
\section{Appendice: le funzioni trigonometriche inverse}
\label{sec:trig-inverse}

Le funzioni trigonometriche $\sen x$, $\cos x$, $\tg x$ non sono monotòne su tutto $\R$ (sono
periodiche) e non esistono le loro funzioni inverse su $\R$: a ogni valore del seno
corrispondono infiniti angoli. Però possiamo restringere a un intervallo limitato
l'insieme in cui prendere in considerazione tali funzioni, in modo che risultino monotòne
nell'insieme considerato. La scelta degli intervalli è convenzionale, ma universalmente
adottata (ed è quella delle funzioni di libreria di tutti i linguaggi di programmazione).

\begin{lemma}\label{lem:monotonia-trig}
La funzione $\sen x$ è strettamente crescente nell'intervallo
$\bigl[-\frac\pi2,\frac\pi2\bigr]$; la funzione $\cos x$ è strettamente decrescente
nell'intervallo $[0,\pi]$; la funzione $\tg x$ è strettamente crescente nell'intervallo
$\bigl(-\frac\pi2,\frac\pi2\bigr)$.
\end{lemma}

\begin{schema}
\Ipotesi{$x_{1}<x_{2}$, entrambi in $\bigl[-\frac\pi2,\frac\pi2\bigr]$ (seno), in $[0,\pi]$
(coseno), in $\bigl(-\frac\pi2,\frac\pi2\bigr)$ (tangente).}
\Tesi{$\sen x_{1}<\sen x_{2}$; \ $\cos x_{1}>\cos x_{2}$; \ $\tg x_{1}<\tg x_{2}$.}
\end{schema}

\begin{strategia}
Non disponiamo ancora del legame tra il segno della derivata e la monotonia (lo
stabiliremo nel prossimo capitolo): studiamo quindi direttamente il segno delle
differenze $\sen x_{2}-\sen x_{1}$, $\cos x_{2}-\cos x_{1}$, $\tg x_{2}-\tg x_{1}$. Le
formule di prostaferesi trasformano le prime due in prodotti di seni e coseni della
semisomma e della semidifferenza, che cadono in intervalli dove il segno di seno e coseno
è noto. Per la tangente si riduce allo stesso denominatore e si riconosce al numeratore il
seno di una differenza.
\end{strategia}

\begin{dimostrazione}
Dalle formule di addizione \eqref{eq:addizione-sen}, \eqref{eq:addizione-cos}, scrivendo
$x_{2}=\frac{x_{1}+x_{2}}2+\frac{x_{2}-x_{1}}2$ e $x_{1}=\frac{x_{1}+x_{2}}2-\frac{x_{2}-x_{1}}2$,
si ottengono le \emph{formule di prostaferesi}
\begin{equation}\label{eq:prostaferesi}
\begin{aligned}
  \sen x_{2}-\sen x_{1}&=2\cos\frac{x_{1}+x_{2}}2\,\sen\frac{x_{2}-x_{1}}2,\\
  \cos x_{2}-\cos x_{1}&=-2\sen\frac{x_{1}+x_{2}}2\,\sen\frac{x_{2}-x_{1}}2 .
\end{aligned}
\end{equation}
Siano $-\frac\pi2\le x_{1}<x_{2}\le\frac\pi2$. Allora $\frac{x_{1}+x_{2}}2\in
\bigl(-\frac\pi2,\frac\pi2\bigr)$, dove il coseno è positivo, e $\frac{x_{2}-x_{1}}2\in
\bigl(0,\frac\pi2\bigr]$, dove il seno è positivo; quindi $\sen x_{2}-\sen x_{1}>0$. Se invece
$0\le x_{1}<x_{2}\le\pi$, allora $\frac{x_{1}+x_{2}}2\in(0,\pi)$ e
$\frac{x_{2}-x_{1}}2\in\bigl(0,\frac\pi2\bigr]$, dove il seno è positivo; quindi
$\cos x_{2}-\cos x_{1}<0$. Infine, se $-\frac\pi2<x_{1}<x_{2}<\frac\pi2$, per la
\eqref{eq:addizione-sen}
\[
  \tg x_{2}-\tg x_{1}=\frac{\sen x_{2}\cos x_{1}-\sen x_{1}\cos x_{2}}{\cos x_{1}\cos x_{2}}
  =\frac{\sen(x_{2}-x_{1})}{\cos x_{1}\cos x_{2}}>0,
\]
perché $x_{2}-x_{1}\in(0,\pi)$ e i due coseni sono positivi.
\end{dimostrazione}

Le formule di prostaferesi sono state usate anche nell'\cref{es:seno-lipschitz} e, nella
forma \eqref{eq:prostaferesi-ab}, nella seconda dimostrazione della
\cref{prop:derivata-sen-cos}.

\subsection{L'arcoseno}

Cominciamo con la funzione $\sen x$. Per il \cref{lem:monotonia-trig} è una funzione
strettamente crescente nell'intervallo $\bigl[-\frac\pi2,\frac\pi2\bigr]$. Consideriamo quindi
$f(x)=\sen x$, con $f:\bigl[-\frac\pi2,\frac\pi2\bigr]\to[-1,1]$. La funzione $f$ è continua e
quindi assume tutti i valori compresi tra il suo minimo ($=-1$) e il suo massimo ($=1$).
Essendo strettamente monotòna, è anche invertibile (\cref{teo:invertibilita}). Pertanto
esiste la funzione inversa
\[
  f^{-1}:[-1,1]\to\Bigl[-\frac\pi2,\frac\pi2\Bigr],
\]
che viene indicata con $f^{-1}(x)=\arcsen x$ (\emph{arcoseno di $x$}). Il nome deriva dal
fatto che, se $y=\arcsen x$, vuol dire che $y$ è uguale alla misura dell'arco, o angolo,
il cui seno è $x$ ($\sen y=x$). Per il \cref{teo:continuita-inversa} l'arcoseno è
continuo, ed è strettamente crescente. Il grafico dell'arcoseno si ottiene
immediatamente dal grafico della funzione seno, per simmetria rispetto alla bisettrice
$y=x$, come nella \cref{fig:arcsen-arccos}.

\begin{figure}[htb]
\centering
\begin{tikzpicture}
\begin{axis}[name=as, width=0.5\textwidth, height=0.5\textwidth, axis lines=middle,
  xmin=-1.9, xmax=1.9, ymin=-1.9, ymax=1.9, axis equal image,
  xtick={-1.5708,-1,1,1.5708}, xticklabels={$-\frac\pi2$,$-1$,$1$,$\frac\pi2$},
  ytick={-1.5708,-1,1,1.5708}, yticklabels={$-\frac\pi2$,$-1$,$1$,$\frac\pi2$},
  tick label style={font=\tiny}, title={\scriptsize $\sen x$ e $\arcsen x$}]
  \addplot[tratteggio, domain=-1.8:1.8] {x};
  \addplot[curva2, domain=-1.5708:1.5708, samples=80] {sin(deg(x))};
  \addplot[curva, domain=-1:1, samples=200] {rad(asin(x))};
\end{axis}
\begin{axis}[at={(as.east)}, anchor=west, xshift=0.8cm, width=0.5\textwidth,
  height=0.5\textwidth, axis lines=middle, xmin=-1.4, xmax=3.5, ymin=-1.4, ymax=3.5,
  axis equal image, xtick={-1,1,1.5708,3.14159}, xticklabels={$-1$,$1$,$\frac\pi2$,$\pi$},
  ytick={-1,1,1.5708,3.14159}, yticklabels={$-1$,$1$,$\frac\pi2$,$\pi$},
  tick label style={font=\tiny}, title={\scriptsize $\cos x$ e $\arccos x$}]
  \addplot[tratteggio, domain=-1.3:3.4] {x};
  \addplot[curva2, domain=0:3.14159, samples=80] {cos(deg(x))};
  \addplot[curva, domain=-1:1, samples=200] {rad(acos(x))};
\end{axis}
\end{tikzpicture}
\caption{A sinistra il seno ristretto a $\bigl[-\frac\pi2,\frac\pi2\bigr]$ (rame) e la sua
inversa $\arcsen x$ (blu), definita su $[-1,1]$; a destra il coseno ristretto a $[0,\pi]$
e la sua inversa $\arccos x$. I grafici delle inverse si ottengono per simmetria rispetto
alla bisettrice $y=x$ (tratteggiata); nei punti $x=\pm1$ hanno tangente verticale.}
\label{fig:arcsen-arccos}
\end{figure}

La funzione arcoseno è quindi definita nell'intervallo \emph{chiuso} $[-1,1]$. Dal teorema
di derivazione delle funzioni inverse si ottiene che la funzione $\arcsen x$ è derivabile
nell'intervallo \emph{aperto} $(-1,1)$, e la derivata vale, posto $y=\arcsen x$,
\begin{equation}\label{eq:derivata-arcsen}
\begin{split}
  D\arcsen x&=\frac1{D\sen y}=\frac1{\cos y}=\frac1{\sqrt{1-\sen^{2}y}}\\
  &=\frac1{\sqrt{1-\sen^{2}(\arcsen x)}}=\frac1{\sqrt{1-x^{2}}};
\end{split}
\end{equation}
sono state utilizzate le relazioni
\begin{equation}\label{eq:relazioni-arcsen}
  \cos y=\sqrt{1-\sen^{2}y}, \qquad \sen(\arcsen x)=x,
\end{equation}
la prima delle quali vale perché $\cos y>0$ per ogni $y\in\bigl(-\frac\pi2,\frac\pi2\bigr)$
(dalla relazione fondamentale, $\cos y=\pm\sqrt{1-\sen^{2}y}$, e il segno è determinato
dalla posizione dell'angolo). La condizione $D\sen y=\cos y\neq0$, richiesta dal
\cref{teo:derivata-inversa}, è appunto verificata per $y$ interno. Nei punti $x=\pm1$, estremi del dominio, si può parlare solo di derivate
unilaterali, e in essi $\cos y=D\sen y$ si annulla ($y=\arcsen(\pm1)=\pm\frac\pi2$). Il
calcolo diretto mostra che il rapporto incrementale tende all'infinito: per $x\to1^{-}$,
posto $y=\arcsen x\to\frac\pi2^{-}$, si ha
\[
  \frac{\arcsen x-\arcsen1}{x-1}=\frac{y-\frac\pi2}{\sen y-1}
  =\frac1{\dfrac{\sen y-\sen\frac\pi2}{y-\frac\pi2}}\to+\infty,
\]
perché il denominatore tende a $D\sen\frac\pi2=\cos\frac\pi2=0$ con valori positivi
(numeratore e denominatore del rapporto $\frac{\sen y-1}{y-\pi/2}$ sono entrambi negativi);
analogamente in $x=-1$. Graficamente, ciò corrisponde al fatto che la funzione
$\arcsen x$ ha retta tangente verticale (unilaterale) se $x=\pm1$.

\subsection{L'arcocoseno}

Passiamo alla funzione $f(x)=\cos x$. Per il \cref{lem:monotonia-trig}, $f:[0,\pi]\to[-1,1]$
è continua e strettamente decrescente; quindi è invertibile in tale intervallo. La
funzione inversa $f^{-1}:[-1,1]\to[0,\pi]$ viene indicata con $f^{-1}(x)=\arccos x$
(\emph{arcocoseno di $x$}). Il grafico si ottiene dal grafico della funzione coseno, come
nella \cref{fig:arcsen-arccos}. La funzione $y=\arccos x$ è definita e continua
nell'intervallo chiuso $[-1,1]$, è derivabile nell'intervallo aperto $(-1,1)$ e la derivata
vale, posto $y=\arccos x\in(0,\pi)$, dove $\sen y>0$,
\begin{equation}\label{eq:derivata-arccos}
\begin{split}
  D\arccos x&=\frac1{D\cos y}=\frac1{-\sen y}=\frac{-1}{\sqrt{1-\cos^{2}y}}\\
  &=\frac{-1}{\sqrt{1-\cos^{2}(\arccos x)}}=\frac{-1}{\sqrt{1-x^{2}}} .
\end{split}
\end{equation}

\begin{osservazione}\label{oss:arcsen-arccos}
Le derivate di $\arcsen x$ e di $\arccos x$ sono opposte; ciò suggerisce che la loro somma
sia costante. In effetti
\begin{equation}\label{eq:arcsen-piu-arccos}
  \arcsen x+\arccos x=\frac\pi2 \qquad\perogni x\in[-1,1].
\end{equation}
Per dimostrarlo, posto $y=\arcsen x\in\bigl[-\frac\pi2,\frac\pi2\bigr]$, l'angolo
$\frac\pi2-y$ appartiene a $[0,\pi]$ e $\cos\bigl(\frac\pi2-y\bigr)=\sen y=x$ (per la
\eqref{eq:addizione-cos}); quindi, per definizione di arcocoseno,
$\arccos x=\frac\pi2-y$. (Nel prossimo capitolo vedremo che una funzione con derivata nulla
su un intervallo è costante, e questa sarà una dimostrazione alternativa.)
\end{osservazione}

\subsection{L'arcotangente}

La più usata funzione trigonometrica inversa è quella relativa alla tangente. La funzione
$f(x)=\tg x$ è continua e strettamente crescente nell'intervallo aperto
$\bigl(-\frac\pi2,\frac\pi2\bigr)$ (\cref{prop:sen-cos-continue,lem:monotonia-trig});
inoltre
\[
  \lim_{x\to\frac\pi2^{-}}\tg x=+\infty, \qquad \lim_{x\to-\frac\pi2^{+}}\tg x=-\infty,
\]
perché $\sen x\to\pm1$ mentre $\cos x\to0$ restando positivo (si veda il riquadro
«$1/0$ non è $+\infty$» del \cref{sec:forme-indeterminate}, con la funzione $\cos x$ al
posto della successione). Dunque, dato un qualsiasi $y\in\R$, esistono $a<b$ in
$\bigl(-\frac\pi2,\frac\pi2\bigr)$ con $\tg a<y<\tg b$, e per il teorema dei valori intermedi
$y$ è un valore della tangente: $\tg$ è una biiezione di $\bigl(-\frac\pi2,\frac\pi2\bigr)$ su
$\R$. La funzione inversa $f^{-1}:\R\to\bigl(-\frac\pi2,\frac\pi2\bigr)$ viene indicata con
$f^{-1}(x)=\operatorname{arctg}x$ (\emph{arcotangente di $x$}), e ha il grafico della
\cref{fig:arctg}. È continua: per ogni $y_{0}\in\R$ la restrizione della tangente a un
intervallo $[a,b]$ come sopra ha inversa continua su $[\tg a,\tg b]$
(\cref{teo:continuita-inversa}), che contiene un intorno di $y_{0}$, e basta applicare
la località del limite.

\begin{figure}[htb]
\centering
\begin{tikzpicture}
\begin{axis}[width=0.9\textwidth, height=5.4cm, axis lines=middle, xmin=-8, xmax=8,
  ymin=-2, ymax=2, xtick={-6,-4,-2,1,2,4,6}, ytick={-1.5708,0.7854,1.5708},
  yticklabels={$-\frac\pi2$,$\frac\pi4$,$\frac\pi2$}, tick label style={font=\scriptsize},
  yticklabel style={fill=white, inner sep=1pt},
  xlabel={$x$}, ylabel={$y$}]
  \addplot[tratteggio, domain=-8:8] {1.5708};
  \addplot[tratteggio, domain=-8:8] {-1.5708};
  \addplot[tratteggio] coordinates {(1,0) (1,0.7854) (0,0.7854)};
  \addplot[curva, domain=-7.8:7.8, samples=200] {rad(atan(x))};
  \node[font=\scriptsize, anchor=south east] at (axis cs:-1.5,0.4) {$y=\operatorname{arctg}x$};
\end{axis}
\end{tikzpicture}
\caption{La funzione $\operatorname{arctg}x$: è definita, continua e strettamente crescente su
tutto $\R$, con asintoti orizzontali $y=\pm\frac\pi2$; si ha $\operatorname{arctg}1=\frac\pi4$.}
\label{fig:arctg}
\end{figure}

La funzione $y=\operatorname{arctg}x$ è definita per ogni $x$ reale. Per i limiti all'infinito
si ha
\begin{equation}\label{eq:arctg-infinito}
  \lim_{x\to+\infty}\operatorname{arctg}x=\frac\pi2, \qquad \lim_{x\to-\infty}\operatorname{arctg}x=-\frac\pi2 .
\end{equation}
Infatti, fissato $\epsilon\in\bigl(0,\frac\pi2\bigr)$, per ogni $x>k=\tg\bigl(\frac\pi2-\epsilon\bigr)$
la monotonia dà $\frac\pi2-\epsilon<\operatorname{arctg}x<\frac\pi2$; il secondo limite segue
per simmetria, essendo $\operatorname{arctg}(-x)=-\operatorname{arctg}x$. È una funzione derivabile per ogni
$x\in\R$, perché $D\tg y=\frac1{\cos^{2}y}\neq0$, e la derivata vale, posto
$y=\operatorname{arctg}x$,
\begin{equation}\label{eq:derivata-arctg}
\begin{split}
  D\operatorname{arctg}x&=\frac1{D\tg y}=\frac1{1/\cos^{2}y}
  =\frac1{(\cos^{2}y+\sen^{2}y)/\cos^{2}y}\\
  &=\frac1{1+\tg^{2}y}=\frac1{1+\tg^{2}(\operatorname{arctg}x)}=\frac1{1+x^{2}} .
\end{split}
\end{equation}

Riassumiamo nella tabella seguente i valori trovati per le derivate.
\begin{center}
\renewcommand{\arraystretch}{1.7}
\begin{tabular}{l l l l}
\toprule
$f(x)$ & dominio & immagine & $f'(x)$\\
\midrule
$\arcsen x$ & $[-1,1]$ & $\bigl[-\frac\pi2,\frac\pi2\bigr]$ & $\dfrac1{\sqrt{1-x^{2}}}$ \ \ $(\abs x<1)$\\
$\arccos x$ & $[-1,1]$ & $[0,\pi]$ & $\dfrac{-1}{\sqrt{1-x^{2}}}$ \ \ $(\abs x<1)$\\
$\operatorname{arctg}x$ & $\R$ & $\bigl(-\frac\pi2,\frac\pi2\bigr)$ & $\dfrac1{1+x^{2}}$\\
\bottomrule
\end{tabular}
\end{center}

\begin{idea}[Gli angoli aerodinamici e la funzione \texttt{atan2}]
In meccanica del volo l'orientamento della velocità rispetto al velivolo si descrive con
due angoli, ricavati dalle componenti $u$ (in avanti), $v$ (laterale), $w$ (verso il basso)
della velocità in assi corpo: l'\emph{angolo d'attacco} $\alpha=\operatorname{arctg}\frac wu$ e
l'\emph{angolo di derapata} $\beta=\arcsen\frac vV$, con $V=\sqrt{u^{2}+v^{2}+w^{2}}$.
L'arcotangente restituisce valori in $\bigl(-\frac\pi2,\frac\pi2\bigr)$ e non distingue,
ad esempio, $(u,w)=(1,1)$ da $(-1,-1)$: in volo normale $u>0$ e non c'è ambiguità, ma in
manovre estreme, in una vite o nei transitori di un modello di simulazione, $u$ può
annullarsi o diventare negativo. Per questo i programmi usano la funzione a due argomenti
$\operatorname{atan2}(w,u)$, che restituisce l'angolo in $(-\pi,\pi]$ del punto $(u,w)$,
tenendo conto del quadrante (e che è definita anche per $u=0$).
\end{idea}

\begin{python}[Angolo d'attacco e angolo di derapata]
import math

# velocità in assi corpo (m/s): u avanti, v laterale, w in basso
casi = [(60.0, 0.0, 5.0), (60.0, 3.0, -4.0), (-2.0, 0.0, 10.0)]
for u, v, w in casi:
    V = math.sqrt(u*u + v*v + w*w)
    alfa_atan = math.degrees(math.atan(w / u))  # in (-90°, 90°)
    alfa = math.degrees(math.atan2(w, u))       # vede il quadrante
    beta = math.degrees(math.asin(v / V))       # derapata
    print(f"u={u:5.1f} v={v:3.1f} w={w:5.1f}: "
          f"atan {alfa_atan:7.2f}°  atan2 {alfa:7.2f}°"
          f"  beta {beta:5.2f}°")
\end{python}
\begin{risultato}
u= 60.0 v=0.0 w=  5.0: atan    4.76°  atan2    4.76°  beta  0.00°
u= 60.0 v=3.0 w= -4.0: atan   -3.81°  atan2   -3.81°  beta  2.86°
u= -2.0 v=0.0 w= 10.0: atan  -78.69°  atan2  101.31°  beta  0.00°
\end{risultato}

Nel terzo caso la velocità, in assi corpo, è diretta quasi del tutto verso il ventre del
velivolo e leggermente all'indietro, come può accadere in una scivolata di coda o in
una manovra post-stallo estrema: l'arcotangente fornisce un
angolo d'attacco di circa $-79^{\circ}$, palesemente sbagliato, mentre $\operatorname{atan2}$
restituisce il valore corretto di circa $101^{\circ}$.

\begin{notastorica}[L'arcotangente e $\pi$]
Le funzioni trigonometriche inverse sono, storicamente, più antiche della loro notazione:
tra il Tre e il Quattrocento il matematico indiano Madhava di Sangamagrama conosceva lo sviluppo in
serie
$\operatorname{arctg}x=x-\frac{x^{3}}3+\frac{x^{5}}5-\dots$, riscoperto in Europa da James Gregory
(1671) e da Leibniz (1674), che ne dedusse, con $x=1$, la celebre formula
$\frac\pi4=1-\frac13+\frac15-\frac17+\cdots$. Formule di questo tipo, combinate con identità
come quella di John Machin (1706), $\frac\pi4=4\operatorname{arctg}\frac15-\operatorname{arctg}\frac1{239}$,
permisero per due secoli di calcolare $\pi$ con centinaia di cifre. La notazione
$\sen^{-1}x$ per l'arcoseno, ancora oggi sui tasti delle calcolatrici, fu introdotta da
John Herschel nel 1813; non va confusa con $1/\sen x$, ed è per evitare l'equivoco che nei
testi italiani si preferisce $\arcsen x$.
\end{notastorica}
